\section{Mục lục và các khối}
\label{sec:muc-luc}

Bảng \textbf{Mục lục} (Outline) ở thanh bên trái liệt kê toàn bộ các mục có
trong tệp đang mở. Bạn có thể bật bảng này bằng biểu tượng danh sách hoặc tổ
hợp phím \phim{Ctrl}+\phim{Shift}+\phim{Y}.

\begin{minipage}[t]{0.28\linewidth}
  \vspace{0pt}
  \khunganh{\linewidth}{muc-luc}
\end{minipage}\hfill
\begin{minipage}[t]{0.66\linewidth}
  \vspace{0pt}
  Khi bấm nút \textbf{Blocks} ở góc trên bên phải bảng, hệ thống sẽ hiển thị
  thêm các \emph{khối} nội dung như hình, bảng, danh sách, công thức, mã nguồn,
  định lý và ghi chú.
  \begin{itemize}
    \item Nhấp vào một dòng để di chuyển ngay tới khối tương ứng.
    \item Rê chuột lên một dòng để xem trước nội dung (như hình ảnh hoặc công
      thức).
    \item Nếu thấy một \textbf{chấm cam}, khối đó đang thiếu thông tin cần
      thiết, chẳng hạn hình chưa có chú thích (caption) hoặc nhãn (label); bạn
      chỉ cần rê chuột lên chấm để xem chi tiết.
    \item Dùng ô lọc \nut{All} để chỉ hiển thị một loại khối nhất định, hoặc lọc
      riêng các khối đang có vấn đề.
  \end{itemize}
\end{minipage}

\subsection{Sửa khối không cần nhớ lệnh}
Khi rê chuột lên một dòng trong Mục lục và bấm vào biểu tượng \textbf{bút chì}
✎, một hộp thoại nhỏ sẽ hiện ra cho phép bạn chỉnh sửa trực quan bằng chuột và
bàn phím. Mọi thay đổi đều được cập nhật thẳng vào mã \LaTeX{}; nếu cần, bạn
có thể nhấn \phim{Ctrl}+\phim{Z} để hoàn tác hoặc \phim{Esc} để đóng hộp thoại.

\begin{meo}
Khi con trỏ đang đặt trong một bảng, hình hoặc công thức, một nút nhỏ
\nut{✎ Edit} sẽ xuất hiện ngay ở đầu khối để bạn mở hộp thoại sửa nhanh. Bạn
cũng có thể dùng phím tắt \phim{Ctrl}+\phim{K} rồi \phim{E}.
\end{meo}

\subsection{Nút ⋯}
Ngay bên cạnh biểu tượng bút chì là nút \textbf{⋯} chứa các thao tác thường
dùng:
\begin{itemize}
  \item \nut{Insert a reference at the cursor}: Chèn trực tiếp tham chiếu tới
    khối (ví dụ \texttt{Figure\textasciitilde\textbackslash ref\{fig:ech\}})
    vào đúng vị trí con trỏ.
  \item \nut{Copy a reference}: Sao chép mã tham chiếu đó để dán vào chỗ khác.
  \item \nut{Duplicate}, \nut{Delete}: Nhân bản hoặc xóa khối.
  \item Riêng với các mục (section), bạn có thêm \nut{Move up} / \nut{Move
    down} để đổi vị trí cả mục, cùng \nut{Promote} / \nut{Demote} để nâng hoặc
    hạ cấp mục.
\end{itemize}

\begin{thuchanh}
Trong sân tập, bạn hãy mở Mục lục, bật chế độ \textbf{Blocks}, sau đó tìm một
khối đang có chấm cam và bổ sung thông tin để chấm cam biến mất.
\end{thuchanh}